% !TEX program = pdflatex
% Jan Andrzejewski — CV
% Kompilacja: pdflatex jan-andrzejewski-cv.tex
% (dla pełnej obsługi „IBM Plex Serif" przełącz na XeLaTeX/LuaLaTeX i odkomentuj sekcję fontspec poniżej)

\documentclass[11pt,a4paper]{article}

\usepackage[T1]{fontenc}
\usepackage[utf8]{inputenc}
\usepackage[polish]{babel}
\usepackage[margin=2cm,top=1.6cm,bottom=1.6cm]{geometry}
\usepackage{enumitem}
\usepackage{hyperref}
\usepackage{titlesec}
\usepackage{xcolor}
\usepackage{parskip}

% --- (opcjonalnie, XeLaTeX/LuaLaTeX) IBM Plex Serif:
% \usepackage{fontspec}
% \setmainfont{IBM Plex Serif}

% --- Kolory zgodne z layoutem Reactive Resume (niebieski akcent)
\definecolor{accent}{HTML}{155DFC}
\definecolor{muted}{HTML}{555555}

\hypersetup{
  colorlinks=true,
  linkcolor=accent,
  urlcolor=accent
}

% --- Nagłówki sekcji: nazwa, kreska pod spodem
\titleformat{\section}
  {\Large\bfseries\color{accent}}{}{0em}{}[{\color{accent}\titlerule[0.6pt]}]
\titlespacing*{\section}{0pt}{1.2em}{0.6em}

\pagestyle{empty}
\setlist[itemize]{leftmargin=1.2em,topsep=2pt,itemsep=1pt,parsep=2pt}

% --- Komenda: tytuł projektu (nazwa + meta po prawej)
\newcommand{\projitem}[2]{%
  \par\medskip
  \noindent{\large\bfseries #1}\hfill{\small\color{muted}\textit{#2}}\par
  \vspace{2pt}
}

% --- Komenda: kategoria umiejętności
\newcommand{\skillrow}[2]{%
  \noindent\textbf{#1} \,\textemdash\, #2\par\vspace{2pt}
}

\begin{document}

% =========================================================
% NAGŁÓWEK
% =========================================================
\begin{center}
  {\Huge\bfseries Jan Andrzejewski}\par
  \vspace{4pt}
  {\large\color{accent} AI-first Developer \,$\cdot$\, Student Automatyki i Robotyki}\par
  \vspace{2pt}
  {\small\color{muted} early adopter LLM, RAG i agentów}\par
  \vspace{6pt}
  {\small
    \href{mailto:janekandrz@gmail.com}{janekandrz@gmail.com}
    \;$\cdot$\; Polska
    \;$\cdot$\; GitHub
    \;$\cdot$\; LinkedIn
  }
\end{center}

% =========================================================
% PODSUMOWANIE
% =========================================================
\section*{Podsumowanie}

\textbf{Inżynier AI / Full-Stack Developer} specjalizujący się w systemach agentycznych
i architekturze systemów produkcyjnych. Łączę solidne fundamenty inżynierskie
(embedded, automatyka) z nowoczesnym stackiem webowym oraz doświadczeniem
w budowaniu pipeline'ów RAG. Early adopter narzędzi AI-assisted development
(od Copilota po Claude Code i Antigravity 2.0). \textbf{W implementacjach kładę
nacisk na świadome decyzje architektoniczne, optymalizację tokenów i wydajność
infrastruktury, co pozwala na realne cięcie kosztów utrzymania systemów AI.}
Student Automatyki i Robotyki realizujący autorskie projekty produktowe end-to-end.

% =========================================================
% DOŚWIADCZENIE
% =========================================================
\section*{Doświadczenie}

\noindent\textbf{PSI Polska} \hfill {\small\color{muted}\href{https://psi.pl}{psi.pl}}\par
\textit{\small(stanowisko i okres do uzupełnienia)}

% =========================================================
% PROJEKTY
% =========================================================
\section*{Projekty}

% --- DocBase
\projitem{DocBase \textemdash{} RAG SaaS „drugi mózg" dla MŚP}%
{2026 \textemdash{} własny projekt produktowy (2-osobowy zespół) \,$\cdot$\, \href{https://docbase.pl}{docbase.pl}}

\textbf{Problem.} W polskich MŚP wiedza operacyjna jest rozproszona po mailach, PDF-ach,
dokumentach i głowach kluczowych pracowników. Onboarding nowej osoby zajmuje tygodnie,
a odejście doświadczonego pracownika to realne ryzyko ciągłości biznesu.

\textbf{Rozwiązanie.} Wielonajemcza platforma SaaS z pipeline'em RAG, która indeksuje
firmowe dokumenty i odpowiada na pytania w języku naturalnym, podając link do źródła.
Zaufanie buduje weryfikowalność odpowiedzi i ścisła izolacja danych każdej firmy
na poziomie bazy.

\textbf{Wartość biznesowa.} Skrócenie onboardingu, oszczędność czasu na szukaniu procedur,
zabezpieczenie wiedzy firmowej. Architektura przygotowana pod opcjonalny wariant
on-premise dla klientów z restrykcyjnym RODO.

\textbf{Decyzje kosztowe i optymalizacje:}
\begin{itemize}
  \item Routing modeli przez \textbf{LiteLLM} (tańszy model do prostych zapytań,
        mocniejszy do analiz), tani model embeddingowy, świadome cięcie tokenów.
  \item Zastąpienie dedykowanej \textbf{bazy wektorowej} zoptymalizowaną tabelą
        \textbf{PostgreSQL}, redukując koszty infrastruktury o rzędy wielkości.
        Przeprowadziłem trade-off między natychmiastowym czasem odpowiedzi
        a kosztami RAM-u, dopasowując architekturę do faktycznych wymagań biznesowych.
\end{itemize}

\textbf{Stack:} Python + \textbf{FastAPI}, \textbf{LlamaIndex},
\textbf{Supabase} (PostgreSQL + Auth + \textbf{RLS}), React + Vite + Tailwind,
\textbf{Docker}, CI/CD na \textbf{GitHub Actions} z deployem landingu na \textbf{AWS S3}.

% --- Tracksy
\projitem{Tracksy \textemdash{} operacyjne SaaS dla salonów beauty}%
{2026 \textemdash{} własny projekt produktowy (2-osobowy zespół)}

\textbf{Problem.} Właściciele i managerowie salonów beauty żyją w Booksy, ale platforma
nie daje im operacyjnego wglądu w to, co realnie kosztuje pieniądze \textemdash{}
puste sloty w grafiku, niewykorzystany czas pracowników, brak szybkich KPI w biegu
między klientami.

\textbf{Rozwiązanie.} Nadbudowa nad Booksy łącząca dwa źródła danych \textemdash{}
historyczny import raportów XLSX (backfill) oraz near real-time strumień rezerwacji,
zmian i odwołań z przychodzących maili przez webhook. Panel zoptymalizowany pod mobile:
KPI i alarmy operacyjne czytelne w kilka sekund.

\textbf{Wartość biznesowa.} Szybka reakcja na luki w grafiku przekłada się bezpośrednio
na przychód, a osobne adresy e-mail per salon pozwalają obsłużyć wiele lokalizacji
bez ingerencji w Booksy.

\textbf{Stack:} \textbf{Next.js} + TypeScript, \textbf{Supabase} (Postgres + Auth + \textbf{RLS}),
\textbf{Supabase Edge Functions} jako odbiornik webhooków,
\textbf{Playwright} do scrapowania panelu Booksy, monorepo na \textbf{Turborepo},
deploy: Vercel + hosted Supabase.

% --- Ball on Beam
\projitem{Stanowisko dydaktyczne Ball on Beam \textemdash{} sterowanie PID/LQR na STM32}%
{2026 \textemdash{} projekt zespołowy, 6. semestr}

\textbf{Problem.} Klasyczny niestabilny obiekt sterowania \textemdash{} kulka swobodnie
tocząca się po pochyłej belce, którą trzeba utrzymać w zadanej pozycji wbrew dynamice
układu. Idealna platforma do praktycznej weryfikacji teorii sterowania na realnym sprzęcie.

\textbf{Rozwiązanie.} Kompletny stos mechatroniczny zbudowany od zera w zespole
\textemdash{} projekt mechaniczny, identyfikacja modelu obiektu, dwa równoległe regulatory
(\textbf{PID} i \textbf{LQR}) z możliwością przełączania w locie oraz desktopowa aplikacja
diagnostyczna z telemetrią i strojeniem nastaw w czasie rzeczywistym.

\textbf{Wartość inżynierska.} Porównanie klasycznego sterowania zwrotnego z optymalnym
sterowaniem w przestrzeni stanów na tym samym obiekcie. Pełny przepływ:
identyfikacja $\rightarrow$ projekt regulatora $\rightarrow$ implementacja na MCU $\rightarrow$ ewaluacja.

\textbf{Stack:} \textbf{STM32} (C, HAL, \textbf{FreeRTOS}), regulatory \textbf{PID} i \textbf{LQR},
czujnik odległości ToF, serwo sterowane PWM, aplikacja desktopowa w \textbf{PySide6}
+ \textbf{pyqtgraph} + \textbf{OpenCV}, identyfikacja modelu w \textbf{MATLAB/Simulink}.

% =========================================================
% EDUKACJA
% =========================================================
\section*{Edukacja}

\noindent\textbf{Studia inżynierskie \textemdash{} Automatyka i Robotyka}
\hfill {\small\color{muted}2023 \textemdash{} obecnie (6. semestr)}\par
Inżynier (w trakcie)

% =========================================================
% UMIEJĘTNOŚCI
% =========================================================
\section*{Umiejętności}

\skillrow{Języki programowania}{Python, TypeScript, C / C++, SQL, MATLAB}
\skillrow{AI / LLM (early adopter)}{RAG (LlamaIndex), OpenAI API, Claude API, Embeddings + pgvector, routing modeli (LiteLLM), Claude Code / Cursor, MCP}
\skillrow{Full-stack web}{Next.js, React, FastAPI, TailwindCSS, Turborepo, Playwright}
\skillrow{Bazy danych \& Cloud}{PostgreSQL / Supabase, Row Level Security, Docker, GitHub Actions (CI/CD), AWS S3, Vercel}
\skillrow{Embedded \& Automatyka}{STM32 + FreeRTOS, Arduino, regulatory PID / LQR, MATLAB / Simulink, OpenCV}

% =========================================================
% JĘZYKI
% =========================================================
\section*{Języki}

\noindent\textbf{Angielski} \textemdash{} biegły w dokumentacji technicznej
i kontakcie z międzynarodowymi klientami.

\end{document}
